\section{精液过敏：一个常被误诊的"怪病"}

前文提到"性交后皮疹、瘙痒可能与精液过敏有关"，本节把这个容易被忽视、又常被误诊为阴道炎或性病的问题讲清楚。男性伴侣也应了解——它是两个人需要共同面对的免疫问题（男性卷有更详细的专章讨论）。

\vspace{0.5em}
\noindent\textbf{它是什么}：精液过敏（医学名"人类精浆超敏反应"，Human Seminal Plasma Hypersensitivity）是女性免疫系统对精液中的蛋白质产生的过敏反应，多为 IgE 介导。真实患病率不高，但因为它"听起来不像病"，很多女性默默忍受多年。

\vspace{0.5em}
\noindent\textbf{典型表现}：

\begin{itemize}
	\item \textbf{局部型}（最常见）：接触后 10--30 分钟内阴道及外阴\textbf{烧灼、瘙痒、肿胀}，可持续数小时至一天以上
	\item \textbf{全身型}（较少但更危险）：全身荨麻疹、面部肿胀、喘息、\textbf{呼吸困难甚至过敏性休克}——多见于对精液高度致敏者
	\item 一个重要的鉴别线索：症状多在\textbf{不使用避孕套的性交后}出现；全程使用避孕套时症状消失或明显减轻
\end{itemize}

\vspace{0.5em}
\noindent\textbf{如何确诊与鉴别}：

\begin{itemize}
	\item \textbf{避孕套试验}是第一步：全程戴套无恙、不戴则发——高度提示精液过敏；但注意乳胶过敏也会在戴套时发作（方向相反），两者都可能被误认为"反复阴道炎"
	\item 与阴道炎、接触性皮炎（对洗涤剂、润滑剂、乳胶过敏）的鉴别需由医生完成：过敏原检测（皮肤点刺、血清特异性 IgE）在专业机构可行
	\item 若每次症状都出现在使用了避孕套或润滑剂之后，则更可能是\textbf{对乳胶或润滑剂成分}过敏，而非精液本身
\end{itemize}

\vspace{0.5em}
\noindent\textbf{管理与治疗}：

\begin{itemize}
	\item \textbf{屏障避孕}：全程避孕套是最直接的手段（乳胶过敏者换聚氨酯套）；备孕需求者需在过敏科/妇科指导下进行
	\item \textbf{药物}：症状出现时可口服抗组胺药；反复发作或有全身反应者，可在性生活前\textbf{预防性}用药（遵医嘱）
	\item \textbf{脱敏治疗}：正规医院可行精液稀释液皮下注射或阴道脱敏方案，多数患者可逐步耐受，\textbf{且需维持规律接触}以防致敏恢复
	\item \textbf{备孕与生育}：精液过敏\textbf{不影响生育能力}；脱敏失败者可通过洗涤精子的人工授精实现怀孕
\end{itemize}

\begin{tcolorbox}[colback=pink!5!white,colframe=red!50!black,title=贴心小叮咛]
性交后反复出现外阴烧灼、肿胀或全身荨麻疹，别再当成"炎症反反复复"硬扛——\textbf{用不用避孕套时症状不一样}这条线索，值得你直接告诉医生。精液过敏不是"私生活不检点"，更不是"心理作用"，它是能查、能治、能正常生育的免疫问题。
\end{tcolorbox}

